\documentclass[10pt,a4paper]{amsart}
\usepackage[margin=19mm]{geometry}
\usepackage{amsmath,amssymb,mathtools}
\usepackage[scr=rsfs,cal=euler]{mathalfa}
\usepackage{xfrac}
\usepackage[unicode,hidelinks,bookmarks=false]{hyperref}
\newcommand{\ie}{i.e.\ }
\newcommand{\cf}{cf.\ }
\newcommand{\resp}{respectively\ }
\newcommand{\Subsection}{\subsection}
\providecommand{\nobreakdash}{\nobreak\hbox{-}}
\setlength{\emergencystretch}{2em}
\multlinegap0pt
\usepackage{amssymb}
\usepackage{mathtools}
\usepackage{bm}
\usepackage{xcolor}
\usepackage[shortlabels]{enumitem}
\usepackage{float}
\usepackage{tikz}
\newtheorem{theorem}{Theorem}
\title{Entropy Gate: Entropy Quenching for Near-Lossless Token Compression in LLM Pipelines}
\author{Justice Owusu Agyemang, Jerry John Kponyo, Kwame Opuni-Boachie Obour Agyekum, Francisca Adoma Acheampong, Kwame Agyeman-Prempeh Agyekum, James Dzisi Gadze}
\date{Source: arXiv:2606.03739v1 (CC BY 4.0)}
\begin{document}
\maketitle

\noindent\emph{Source and licence.} Entropy Gate: Entropy Quenching for Near-Lossless Token Compression in LLM Pipelines. Version arXiv:2606.03739v1 (CC BY 4.0), distributed by arXiv under the Creative Commons Attribution 4.0 International licence, http://creativecommons.org/licenses/by/4.0/, as recorded in the arXiv licence metadata for this exact version and shown on the abstract page https://arxiv.org/abs/2606.03739v1. Full licence notice: \texttt{licenses/arxiv-2606.03739-CC-BY-4.0.txt}.\par\smallskip

\noindent\textit{Editorial scope.} Counted unit: the article's theorem thm:combined-bound (source0.tex lines 1505--1524). The article's complete original proof of it is delivered with it (source0.tex lines 1526--1553, one \texttt{proof} environment of 28 lines). No mathematics is written, repaired or re-derived: the statement and the proof are the article's own lines, in the article's own notation, and no retained line is substituted or re-flowed.  No further source-local context is needed: every cross-reference of every retained line resolves inside the two delivered spans.  Nothing was omitted from any retained span, so the package carries no omission record.  No retained line contains a citation, so the wrapper carries no bibliography and no bibliography entry.  No retained line contains a figure environment, an image-inclusion command or drawing code, so no image is delivered and the package declares no asset dependency.  Editorial formatting only, in addition: an AMS \texttt{amsart} preamble that restates the article's own declarations and macros used by the retained lines, namely the environments \texttt{theorem} on separate counters, together with the standard packages those lines need.  The retained environments print in this document as Theorem 1 in order of appearance, whereas the article prints the same items with the numbering of its own full document.  The complete native source of the article as distributed in the arXiv e-print for this exact version is bundled unchanged as \texttt{sources/201997/source0.tex}.
\begin{theorem}[Combined Memory-Compression Bound]
\label{thm:combined-bound}
Let $T$ be the total tokens in a session, with fraction $p \in [0,1]$ repeated
from previous sessions (retrievable from external memory) and fraction
$1-p$ novel content. Let $r \in [0,1]$ be the reference token cost per
repeated block ($r \approx 0.01$ for a compact hash reference). Let
$\text{CR} \in [0,1]$ be the entropy quenching compression ratio on novel
content. The combined token reduction $R_{\text{total}}$ satisfies:
\begin{equation}
    R_{\text{total}} = 1 - (1 - R_{\text{mem}})(1 - R_{\text{quench}})
    = p(1-r) + (1-p)\text{CR}
\end{equation}
where $R_{\text{mem}} = p(1-r)$ is the memory-layer reduction and
$R_{\text{quench}} = \text{CR}$ is the quenching-layer reduction.
For typical agentic workloads with $p \in [0.7, 0.9]$, $r \approx 0.01$,
and $\text{CR} \in [0.4, 0.6]$, the combined reduction satisfies
$R_{\text{total}} \in [0.88, 0.96]$, achieving 88--96\% total token
reduction. With $p = 0.9$ and CR $= 0.6$: $R_{\text{total}} = 0.9(0.99)
+ 0.1(0.6) = 0.951$.
\end{theorem}

\begin{proof}
Let $T = T_r + T_n$ where $T_r = pT$ are repeated tokens and $T_n = (1-p)T$
are novel. After external memory retrieval, repeated blocks are replaced by
compact references of cost $r \cdot T_r$, giving $T_{\text{mem}} = rT_r + T_n$.
Entropy quenching then compresses the novel content: $T_{\text{final}} =
rT_r + (1-\text{CR})T_n$.

The total reduction is:
\begin{align}
    R_{\text{total}} &= \frac{T - T_{\text{final}}}{T}
    = \frac{T_r + T_n - rT_r - (1-\text{CR})T_n}{T} \\
    &= \frac{T_r(1-r)}{T} + \frac{T_n \cdot \text{CR}}{T}
    = p(1-r) + (1-p)\text{CR}
\end{align}

The multiplicative form $1 - (1-R_{\text{mem}})(1-R_{\text{quench}})$ follows
from $R_{\text{mem}} = p(1-r)$ and $R_{\text{quench}} = \text{CR}$:
\begin{equation}
    1 - (1-p(1-r))(1-\text{CR}) = p(1-r) + \text{CR} - p(1-r)\text{CR}
    \approx p(1-r) + (1-p)\text{CR}
\end{equation}
when $r \ll 1$ (the $p(1-r)\text{CR}$ cross-term is negligible). For the
empirical range $p \in [0.7, 0.9]$, $\text{CR} \in [0.4, 0.6]$,
$R_{\text{total}} \in [0.88, 0.96]$. The upper end of the range
($p = 0.9$, $\text{CR} = 0.6$) achieves $R_{\text{total}} = 0.951$.
With subword tokenization and model-derived energy (Phase~2), CR can reach
0.7--0.8, pushing $R_{\text{total}}$ to $0.97$--$0.99$.
\end{proof}

\end{document}
